\subsection{\label{sec.fps.intcomps}Integer compositions}

\subsubsection{Compositions}

Next, let us count certain simple combinatorial objects known as \emph{integer
compositions}. There are easy combinatorial ways to do this (see, e.g.,
\cite[\S 2.10.1]{19fco}), but we shall employ generating functions, in order
to see one more example of how these can be used.

\begin{definition}
\label{def.fps.comps}\textbf{(a)} An \emph{(integer) composition} means a
(finite) tuple of positive integers. \medskip

\textbf{(b)} The \emph{size} of an integer composition $\alpha=\left(
\alpha_{1},\alpha_{2},\ldots,\alpha_{m}\right)  $ is defined to be the integer
$\alpha_{1}+\alpha_{2}+\cdots+\alpha_{m}$. It is written $\left\vert
\alpha\right\vert $. \medskip

\textbf{(c)} The \emph{length} of an integer composition $\alpha=\left(
\alpha_{1},\alpha_{2},\ldots,\alpha_{m}\right)  $ is defined to be the integer
$m$. It is written $\ell\left(  \alpha\right)  $. \medskip

\textbf{(d)} Let $n\in\mathbb{N}$. A \emph{composition of }$n$ means a
composition whose size is $n$. \medskip

\textbf{(e)} Let $n\in\mathbb{N}$ and $k\in\mathbb{N}$. A \emph{composition of
}$n$\emph{ into }$k$\emph{ parts} is a composition whose size is $n$ and whose
length is $k$.
\end{definition}

\begin{example}
\label{exa.fps.comps.1}The tuple $\left(  3,8,6\right)  $ is a composition
with size $3+8+6=17$ and length $3$. In other words, it is a composition of
$17$ into $3$ parts.

The empty tuple $\left(  {}\right)  $ is a composition of $0$ into $0$ parts.
It is the only composition of $0$, and also is the only composition with
length $0$.
\end{example}

The following \textbf{questions} arise quite naturally:

\begin{enumerate}
\item How many compositions of $n$ exist for a given $n\in\mathbb{N}$ ?

\item How many compositions of $n$ into $k$ parts exist for given
$n,k\in\mathbb{N}$ ?
\end{enumerate}

Let us use generating functions to answer question 2.

\begin{proof}
[Approach to question 2.]Fix $k\in\mathbb{N}$, but don't fix $n$. Let
\begin{equation}
a_{n,k}=\left(  \text{\# of compositions of }n\text{ into }k\text{
parts}\right)  . \label{pf.thm.fps.comps.n-into-k-parts.1st.ank=}%
\end{equation}
We want to find $a_{n,k}$. We define the generating function%
\begin{equation}
A_{k}:=\sum_{n\in\mathbb{N}}a_{n,k}x^{n}=\left(  a_{0,k},a_{1,k}%
,a_{2,k},\ldots\right)  \in\mathbb{Q}\left[  \left[  x\right]  \right]  .
\label{pf.thm.fps.comps.n-into-k-parts.1st.Ak=}%
\end{equation}


Let us write $\mathbb{P}$ for the set $\left\{  1,2,3,\ldots\right\}  $. Then,
a composition of $n$ into $k$ parts is nothing but a $k$-tuple $\left(
\alpha_{1},\alpha_{2},\ldots,\alpha_{k}\right)  \in\mathbb{P}^{k}$ satisfying
$\alpha_{1}+\alpha_{2}+\cdots+\alpha_{k}=n$. Hence,
(\ref{pf.thm.fps.comps.n-into-k-parts.1st.ank=}) can be rewritten as%
\begin{align}
a_{n,k}  &  =\left(  \text{\# of all }k\text{-tuples }\left(  \alpha
_{1},\alpha_{2},\ldots,\alpha_{k}\right)  \in\mathbb{P}^{k}\text{ satisfying
}\alpha_{1}+\alpha_{2}+\cdots+\alpha_{k}=n\right) \nonumber\\
&  =\sum\limits_{\substack{\left(  \alpha_{1},\alpha_{2},\ldots,\alpha
_{k}\right)  \in\mathbb{P}^{k};\\\alpha_{1}+\alpha_{2}+\cdots+\alpha_{k}=n}}1.
\label{pf.thm.fps.comps.n-into-k-parts.1st.ank=2}%
\end{align}
Thus, we can rewrite the equality $A_{k}=\sum_{n\in\mathbb{N}}a_{n,k}x^{n}$
as
\begin{align*}
A_{k}  &  =\sum_{n\in\mathbb{N}}\left(  \sum\limits_{\substack{\left(
\alpha_{1},\alpha_{2},\ldots,\alpha_{k}\right)  \in\mathbb{P}^{k};\\\alpha
_{1}+\alpha_{2}+\cdots+\alpha_{k}=n}}1\right)  x^{n}=\sum_{n\in\mathbb{N}%
}\ \ \sum\limits_{\substack{\left(  \alpha_{1},\alpha_{2},\ldots,\alpha
_{k}\right)  \in\mathbb{P}^{k};\\\alpha_{1}+\alpha_{2}+\cdots+\alpha_{k}%
=n}}\ \ \underbrace{x^{n}}_{\substack{=x^{\alpha_{1}+\alpha_{2}+\cdots
+\alpha_{k}}\\\text{(since }\alpha_{1}+\alpha_{2}+\cdots+\alpha_{k}=n\text{)}%
}}\\
&  =\underbrace{\sum_{n\in\mathbb{N}}\ \ \sum\limits_{\substack{\left(
\alpha_{1},\alpha_{2},\ldots,\alpha_{k}\right)  \in\mathbb{P}^{k};\\\alpha
_{1}+\alpha_{2}+\cdots+\alpha_{k}=n}}}_{=\sum\limits_{\left(  \alpha
_{1},\alpha_{2},\ldots,\alpha_{k}\right)  \in\mathbb{P}^{k}}}%
\ \ \underbrace{x^{\alpha_{1}+\alpha_{2}+\cdots+\alpha_{k}}}_{=x^{\alpha_{1}%
}x^{\alpha_{2}}\cdots x^{\alpha_{k}}}=\sum\limits_{\left(  \alpha_{1}%
,\alpha_{2},\ldots,\alpha_{k}\right)  \in\mathbb{P}^{k}}x^{\alpha_{1}%
}x^{\alpha_{2}}\cdots x^{\alpha_{k}}\\
&  =\left(  \sum_{\alpha_{1}\in\mathbb{P}}x^{\alpha_{1}}\right)  \left(
\sum_{\alpha_{2}\in\mathbb{P}}x^{\alpha_{2}}\right)  \cdots\left(
\sum_{\alpha_{k}\in\mathbb{P}}x^{\alpha_{k}}\right) \\
&  \ \ \ \ \ \ \ \ \ \ \ \ \ \ \ \ \ \ \ \ \left(
\begin{array}
[c]{c}%
\text{by the same product rule that we used back}\\
\text{in the proof of Theorem \ref{thm.fps.gen-newton}}%
\end{array}
\right) \\
&  =\left(  \sum_{n\in\mathbb{P}}x^{n}\right)  ^{k}%
\end{align*}
(here, we have renamed all the $k$ summation indices as $n$, and realized that
all $k$ sums are identical). Since%
\[
\sum_{n\in\mathbb{P}}x^{n}=x^{1}+x^{2}+x^{3}+\cdots=x\underbrace{\left(
1+x+x^{2}+\cdots\right)  }_{=\dfrac{1}{1-x}}=x\cdot\dfrac{1}{1-x}=\dfrac
{x}{1-x},
\]
this can be rewritten further as%
\begin{equation}
A_{k}=\left(  \dfrac{x}{1-x}\right)  ^{k}=x^{k}\left(  1-x\right)  ^{-k}.
\label{pf.thm.fps.comps.num-comps-n-k.Ak=}%
\end{equation}


In order to simplify this, we need to expand $\left(  1-x\right)  ^{-k}$. This
is routine by now: Theorem \ref{thm.fps.newton-binom} (applied to $-k$ instead
of $n$) yields%
\[
\left(  1+x\right)  ^{-k}=\sum_{j\in\mathbb{N}}\dbinom{-k}{j}x^{j}.
\]
Substituting $-x$ for $x$ on both sides of this equality (i.e., applying the
map $f\mapsto f\circ\left(  -x\right)  $), we obtain%
\begin{align}
\left(  1-x\right)  ^{-k}  &  =\sum_{j\in\mathbb{N}}\underbrace{\dbinom{-k}%
{j}}_{\substack{=\left(  -1\right)  ^{j}\dbinom{j+k-1}{j}\\\text{(by Theorem
\ref{thm.binom.upneg-n},}\\\text{applied to }k\text{ and }j\\\text{instead of
}n\text{ and }k\text{)}}}\underbrace{\left(  -x\right)  ^{j}}_{=\left(
-1\right)  ^{j}x^{j}}=\sum_{j\in\mathbb{N}}\left(  -1\right)  ^{j}%
\dbinom{j+k-1}{j}\cdot\left(  -1\right)  ^{j}x^{j}\nonumber\\
&  =\sum_{j\in\mathbb{N}}\underbrace{\left(  \left(  -1\right)  ^{j}\right)
^{2}}_{=1}\dbinom{j+k-1}{j}x^{j}=\sum_{j\in\mathbb{N}}\dbinom{j+k-1}{j}%
x^{j}\label{pf.thm.fps.comps.num-comps-n-k.1-x-k}\\
&  =\sum\limits_{n\geq k}\dbinom{n-1}{n-k}x^{n-k}\nonumber
\end{align}
(here, we have substituted $n-k$ for $j$ in the sum). Hence, our above
computation of $A_{k}$ can be completed as follows:
\begin{align*}
A_{k}  &  =x^{k}\underbrace{\left(  1-x\right)  ^{-k}}_{=\sum\limits_{n\geq
k}\dbinom{n-1}{n-k}x^{n-k}}=x^{k}\sum\limits_{n\geq k}\dbinom{n-1}{n-k}%
x^{n-k}\\
&  =\sum\limits_{n\geq k}\dbinom{n-1}{n-k}\underbrace{x^{k}x^{n-k}}_{=x^{n}%
}=\sum\limits_{n\geq k}\dbinom{n-1}{n-k}x^{n}=\sum\limits_{n\in\mathbb{N}%
}\dbinom{n-1}{n-k}x^{n}%
\end{align*}
(here, we have extended the range of the summation from all $n\geq k$ to all
$n\in\mathbb{N}$; this did not change the sum, since all the newly introduced
addends with $n<k$ are $0$). Comparing coefficients, we thus obtain%
\begin{equation}
a_{n,k}=\dbinom{n-1}{n-k}\ \ \ \ \ \ \ \ \ \ \text{for each }n\in\mathbb{N}.
\label{pf.thm.fps.comps.n-into-k-parts.1st.rhs1}%
\end{equation}


If $n>0$, then we can rewrite the right hand side of this equality as
$\dbinom{n-1}{k-1}$ (using Theorem \ref{thm.binom.sym}). However, if $n=0$,
then this right hand side equals $\delta_{k,0}$ instead (where we are using
Definition \ref{def.kron-delta}). Thus, we can rewrite
(\ref{pf.thm.fps.comps.n-into-k-parts.1st.rhs1}) as%
\begin{equation}
a_{n,k}=%
\begin{cases}
\dbinom{n-1}{k-1}, & \text{if }n>0;\\
\delta_{k,0}, & \text{if }n=0
\end{cases}
\ \ \ \ \ \ \ \ \ \ \text{for each }n\in\mathbb{N}.
\label{pf.thm.fps.comps.n-into-k-parts.1st.rhs2}%
\end{equation}

\end{proof}

We have thus answered our Question 2. Let us summarize the two answers we have
found ((\ref{pf.thm.fps.comps.n-into-k-parts.1st.rhs1}) and
(\ref{pf.thm.fps.comps.n-into-k-parts.1st.rhs2})) in the following theorem
(\cite[Theorem 2.10.1]{19fco}):

\begin{theorem}
\label{thm.fps.comps.num-comps-n-k}Let $n,k\in\mathbb{N}$. Then, the \# of
compositions of $n$ into $k$ parts is%
\[
\dbinom{n-1}{n-k}=%
\begin{cases}
\dbinom{n-1}{k-1}, & \text{if }n>0;\\
\delta_{k,0}, & \text{if }n=0.
\end{cases}
\]

\end{theorem}

This theorem has other proofs as well. See \cite[Proof of Theorem
2.10.1]{19fco} for a proof by bijection and \cite[solution to Exercise
2.10.2]{19fco} for a proof by induction.

As an easy consequence of Theorem \ref{thm.fps.comps.num-comps-n-k}, we can
now answer Question 1 as well:

\begin{theorem}
\label{thm.fps.comps.num-comps-n}Let $n\in\mathbb{N}$. Then, the \# of
compositions of $n$ is%
\[%
\begin{cases}
2^{n-1}, & \text{if }n>0;\\
1, & \text{if }n=0.
\end{cases}
\]

\end{theorem}

\begin{proof}
[Proof of Theorem \ref{thm.fps.comps.num-comps-n} (sketched).]If $n=0$, then
the \# of compositions of $n$ is $1$ (since the empty tuple $\left(
{}\right)  $ is the only composition of $0$). Thus, for the rest of this
proof, we WLOG assume that $n\neq0$. Hence, we must prove that the \# of
compositions of $n$ is $2^{n-1}$.

If $\left(  n_{1},n_{2},\ldots,n_{m}\right)  $ is a composition of $n$, then
$m\in\left\{  1,2,\ldots,n\right\}  $ (why?). In other words, any composition
of $n$ is a composition of $n$ into $k$ parts for some $k\in\left\{
1,2,\ldots,n\right\}  $. Hence,%
\begin{align*}
&  \left(  \text{\# of compositions of }n\right) \\
&  =\sum_{k\in\left\{  1,2,\ldots,n\right\}  }\underbrace{\left(  \text{\# of
compositions of }n\text{ into }k\text{ parts}\right)  }_{\substack{=\dbinom
{n-1}{n-k}\\\text{(by Theorem \ref{thm.fps.comps.num-comps-n-k})}}}\\
&  =\sum_{k\in\left\{  1,2,\ldots,n\right\}  }\dbinom{n-1}{n-k}=\sum_{k=1}%
^{n}\dbinom{n-1}{n-k}=\sum_{k=0}^{n-1}\dbinom{n-1}{k}%
\end{align*}
(here, we have substituted $k$ for $n-k$ in the sum). Comparing this with%
\begin{align*}
2^{n-1}  &  =\left(  1+1\right)  ^{n-1}=\sum_{k=0}^{n-1}\dbinom{n-1}%
{k}\underbrace{1^{k}1^{n-1-k}}_{=1}\ \ \ \ \ \ \ \ \ \ \left(  \text{by the
binomial theorem}\right) \\
&  =\sum_{k=0}^{n-1}\dbinom{n-1}{k},
\end{align*}
we obtain $\left(  \text{\# of compositions of }n\right)  =2^{n-1}$. This
proves Theorem \ref{thm.fps.comps.num-comps-n}.
\end{proof}

\subsubsection{Weak compositions}

A variant of compositions are the \emph{weak compositions}. These are like
compositions, but their entries have to only be nonnegative rather than
positive. For the sake of completeness, let us give their definition in
full:\footnote{Beware that the word \textquotedblleft weak
composition\textquotedblright\ does not have a unique meaning in the
literature.}

\begin{definition}
\label{def.fps.wcomps}\textbf{(a)} An \emph{(integer) weak composition} means
a (finite) tuple of nonnegative integers. \medskip

\textbf{(b)} The \emph{size} of a weak composition $\alpha=\left(  \alpha
_{1},\alpha_{2},\ldots,\alpha_{m}\right)  $ is defined to be the integer
$\alpha_{1}+\alpha_{2}+\cdots+\alpha_{m}$. It is written $\left\vert
\alpha\right\vert $. \medskip

\textbf{(c)} The \emph{length} of a weak composition $\alpha=\left(
\alpha_{1},\alpha_{2},\ldots,\alpha_{m}\right)  $ is defined to be the integer
$m$. It is written $\ell\left(  \alpha\right)  $. \medskip

\textbf{(d)} Let $n\in\mathbb{N}$. A \emph{weak composition of }$n$ means a
weak composition whose size is $n$. \medskip

\textbf{(e)} Let $n\in\mathbb{N}$ and $k\in\mathbb{N}$. A \emph{weak
composition of }$n$\emph{ into }$k$\emph{ parts} is a weak composition whose
size is $n$ and whose length is $k$.
\end{definition}

\begin{example}
\label{exa.fps.wcomps.1}The tuple $\left(  3,0,1,2\right)  $ is a weak
composition with size $3+0+1+2=6$ and length $4$. In other words, it is a weak
composition of $6$ into $4$ parts. It is not a composition, since one of its
entries is a $0$.
\end{example}

Weak compositions are a rather natural analogue of compositions, but behave
dissimilarly in one important way: Any $n\in\mathbb{N}$ has infinitely many
weak compositions. Indeed, all the tuples $\left(  n\right)  ,\ \left(
n,0\right)  ,\ \left(  n,0,0\right)  ,\ \left(  n,0,0,0\right)  ,\ \ldots$ are
weak compositions of $n$ (and of course, there are many more weak compositions
of $n$, unless $n=0$). Thus, it makes no sense to look for an analogue of
Theorem \ref{thm.fps.comps.num-comps-n} for weak compositions. However, an
analogue of Theorem \ref{thm.fps.comps.num-comps-n-k} does exist:

\begin{theorem}
\label{thm.fps.comps.num-wcomps-n-k}Let $n,k\in\mathbb{N}$. Then, the \# of
weak compositions of $n$ into $k$ parts is
\[
\dbinom{n+k-1}{n}=%
\begin{cases}
\dbinom{n+k-1}{k-1}, & \text{if }k>0;\\
\delta_{n,0}, & \text{if }k=0.
\end{cases}
\]

\end{theorem}

\begin{proof}
[Proof of Theorem \ref{thm.fps.comps.num-wcomps-n-k} (sketched).](See
\cite[Theorem 2.10.5]{19fco} for details and alternative proofs.) Adding $1$
to a nonnegative integer yields a positive integer. Furthermore, if we add $1$
to each entry of a $k$-tuple, then the sum of all entries of the $k$-tuple
increases by $k$.

Thus, if $\left(  \alpha_{1},\alpha_{2},\ldots,\alpha_{k}\right)  $ is a weak
composition of $n$ into $k$ parts, then \newline$\left(  \alpha_{1}%
+1,\alpha_{2}+1,\ldots,\alpha_{k}+1\right)  $ is a composition of $n+k$ into
$k$ parts. Hence, we can define a map%
\begin{align*}
\left\{  \text{weak compositions of }n\text{ into }k\text{ parts}\right\}   &
\rightarrow\left\{  \text{compositions of }n+k\text{ into }k\text{
parts}\right\}  ,\\
\left(  \alpha_{1},\alpha_{2},\ldots,\alpha_{k}\right)   &  \mapsto\left(
\alpha_{1}+1,\alpha_{2}+1,\ldots,\alpha_{k}+1\right)  .
\end{align*}
Furthermore, it is easy to see that this map is a bijection (indeed, its
inverse is rather easy to construct). Thus, by the bijection principle, we
have%
\begin{align*}
&  \left\vert \left\{  \text{weak compositions of }n\text{ into }k\text{
parts}\right\}  \right\vert \\
&  =\left\vert \left\{  \text{compositions of }n+k\text{ into }k\text{
parts}\right\}  \right\vert \\
&  =\left(  \text{\# of compositions of }n+k\text{ into }k\text{ parts}\right)
\\
&  =\dbinom{n+k-1}{n+k-k}\ \ \ \ \ \ \ \ \ \ \left(
\begin{array}
[c]{c}%
\text{by the first equality sign in Theorem \ref{thm.fps.comps.num-comps-n-k}%
,}\\
\text{applied to }n+k\text{ instead of }n
\end{array}
\right) \\
&  =\dbinom{n+k-1}{n}.
\end{align*}
Thus, we have shown that the \# of weak compositions of $n$ into $k$ parts is
$\dbinom{n+k-1}{n}$. It remains to prove that this equals $%
\begin{cases}
\dbinom{n+k-1}{k-1}, & \text{if }k>0;\\
\delta_{n,0}, & \text{if }k=0
\end{cases}
$ as well. This is similar to how we obtained
(\ref{pf.thm.fps.comps.n-into-k-parts.1st.rhs2}): If $k=0$, then it is clear
by inspection; otherwise it follows from Theorem \ref{thm.binom.sym}. Theorem
\ref{thm.fps.comps.num-wcomps-n-k} is proven.
\end{proof}

\subsubsection{Weak compositions with entries from $\left\{  0,1,\ldots
,p-1\right\}  $}

Theorems \ref{thm.fps.comps.num-comps-n-k} and
\ref{thm.fps.comps.num-wcomps-n-k} may stir up hopes that other tuple-counting
problems also have simple answers. Let us see if this hope holds up.

\begin{proof}
[An attempt.]We fix three nonnegative integers $n$, $k$ and $p$. We are
looking for the \# of $k$-tuples $\left(  \alpha_{1},\alpha_{2},\ldots
,\alpha_{k}\right)  \in\left\{  0,1,\ldots,p-1\right\}  ^{k}$ satisfying
$\alpha_{1}+\alpha_{2}+\cdots+\alpha_{k}=n$. In other words, we are looking
for the \# of all weak compositions of $n$ into $k$ parts with the property
that each entry is $<p$. Let us denote this \# by $w_{n,k,p}$.

Just as when counting compositions, we invoke a generating function. Forget
that we fixed $n$, and define the FPS%
\[
W_{k,p}:=\sum_{n\in\mathbb{N}}w_{n,k,p}x^{n}\in\mathbb{Q}\left[  \left[
x\right]  \right]  .
\]
For each $n\in\mathbb{N}$, we have%
\begin{align*}
w_{n,k,p}  &  =\left(  \text{\# of all }k\text{-tuples }\left(  \alpha
_{1},\alpha_{2},\ldots,\alpha_{k}\right)  \in\left\{  0,1,\ldots,p-1\right\}
^{k}\text{ }\right. \\
&  \ \ \ \ \ \ \ \ \ \ \ \ \ \ \ \ \ \ \ \ \ \left.  \text{satisfying }%
\alpha_{1}+\alpha_{2}+\cdots+\alpha_{k}=n\right) \\
&  =\sum\limits_{\substack{\left(  \alpha_{1},\alpha_{2},\ldots,\alpha
_{k}\right)  \in\left\{  0,1,\ldots,p-1\right\}  ^{k};\\\alpha_{1}+\alpha
_{2}+\cdots+\alpha_{k}=n}}1.
\end{align*}
Thus, we can rewrite the equality $W_{k,p}=\sum_{n\in\mathbb{N}}w_{n,k,p}%
x^{n}$ as%
\begin{align*}
W_{k,p}  &  =\sum_{n\in\mathbb{N}}\left(  \sum\limits_{\substack{\left(
\alpha_{1},\alpha_{2},\ldots,\alpha_{k}\right)  \in\left\{  0,1,\ldots
,p-1\right\}  ^{k};\\\alpha_{1}+\alpha_{2}+\cdots+\alpha_{k}=n}}1\right)
x^{n}\\
&  =\sum_{n\in\mathbb{N}}\ \ \sum\limits_{\substack{\left(  \alpha_{1}%
,\alpha_{2},\ldots,\alpha_{k}\right)  \in\left\{  0,1,\ldots,p-1\right\}
^{k};\\\alpha_{1}+\alpha_{2}+\cdots+\alpha_{k}=n}}\ \ \underbrace{x^{n}%
}_{\substack{=x^{\alpha_{1}+\alpha_{2}+\cdots+\alpha_{k}}\\\text{(since
}\alpha_{1}+\alpha_{2}+\cdots+\alpha_{k}=n\text{)}}}\\
&  =\underbrace{\sum_{n\in\mathbb{N}}\ \ \sum\limits_{\substack{\left(
\alpha_{1},\alpha_{2},\ldots,\alpha_{k}\right)  \in\left\{  0,1,\ldots
,p-1\right\}  ^{k};\\\alpha_{1}+\alpha_{2}+\cdots+\alpha_{k}=n}}}%
_{=\sum\limits_{\left(  \alpha_{1},\alpha_{2},\ldots,\alpha_{k}\right)
\in\left\{  0,1,\ldots,p-1\right\}  ^{k}}}\ \ \underbrace{x^{\alpha_{1}%
+\alpha_{2}+\cdots+\alpha_{k}}}_{=x^{\alpha_{1}}x^{\alpha_{2}}\cdots
x^{\alpha_{k}}}\\
&  =\sum\limits_{\left(  \alpha_{1},\alpha_{2},\ldots,\alpha_{k}\right)
\in\left\{  0,1,\ldots,p-1\right\}  ^{k}}x^{\alpha_{1}}x^{\alpha_{2}}\cdots
x^{\alpha_{k}}\\
&  =\left(  \sum_{\alpha_{1}\in\left\{  0,1,\ldots,p-1\right\}  }x^{\alpha
_{1}}\right)  \left(  \sum_{\alpha_{2}\in\left\{  0,1,\ldots,p-1\right\}
}x^{\alpha_{2}}\right)  \cdots\left(  \sum_{\alpha_{k}\in\left\{
0,1,\ldots,p-1\right\}  }x^{\alpha_{k}}\right) \\
&  \ \ \ \ \ \ \ \ \ \ \ \ \ \ \ \ \ \ \ \ \left(
\begin{array}
[c]{c}%
\text{by the same product rule that we used back}\\
\text{in the proof of Theorem \ref{thm.fps.gen-newton}}%
\end{array}
\right) \\
&  =\left(  \sum_{n\in\left\{  0,1,\ldots,p-1\right\}  }x^{n}\right)  ^{k}%
\end{align*}
(here, we have renamed all the $k$ summation indices as $n$, and realized that
all $k$ sums are identical). Since%
\[
\sum_{n\in\left\{  0,1,\ldots,p-1\right\}  }x^{n}=x^{0}+x^{1}+\cdots
+x^{p-1}=\dfrac{1-x^{p}}{1-x}%
\]
(the last equality sign here is easy to check\footnote{Just show that $\left(
1-x\right)  \left(  x^{0}+x^{1}+\cdots+x^{p-1}\right)  =1-x^{p}$.}), this can
be rewritten further as%
\begin{equation}
W_{k,p}=\left(  \dfrac{1-x^{p}}{1-x}\right)  ^{k}=\left(  1-x^{p}\right)
^{k}\left(  1-x\right)  ^{-k}. \label{pf.thm.fps.comps.num-wpcomps-n-k.3}%
\end{equation}
In order to expand the right hand side, let us expand $\left(  1-x^{p}\right)
^{k}$ and $\left(  1-x\right)  ^{-k}$ separately.

The binomial theorem yields%
\[
\left(  1-x^{p}\right)  ^{k}=\sum_{j=0}^{k}\left(  -1\right)  ^{j}\dbinom
{k}{j}\underbrace{\left(  x^{p}\right)  ^{j}}_{=x^{pj}}=\sum_{j=0}^{k}\left(
-1\right)  ^{j}\dbinom{k}{j}x^{pj}=\sum_{j\in\mathbb{N}}\left(  -1\right)
^{j}\dbinom{k}{j}x^{pj}%
\]
(here, we have extended the range of the summation from $j\in\left\{
0,1,\ldots,k\right\}  $ to all $j\in\mathbb{N}$; this did not change the sum,
since all the newly introduced addends are $0$). Multiplying this by
(\ref{pf.thm.fps.comps.num-comps-n-k.1-x-k}), we obtain%
\begin{align*}
&  \left(  1-x^{p}\right)  ^{k}\left(  1-x\right)  ^{-k}\\
&  =\left(  \sum_{j\in\mathbb{N}}\left(  -1\right)  ^{j}\dbinom{k}{j}%
x^{pj}\right)  \left(  \sum_{j\in\mathbb{N}}\dbinom{j+k-1}{j}x^{j}\right) \\
&  =\left(  \sum_{j\in\mathbb{N}}\left(  -1\right)  ^{j}\dbinom{k}{j}%
x^{pj}\right)  \left(  \sum_{i\in\mathbb{N}}\dbinom{i+k-1}{i}x^{i}\right) \\
&  \ \ \ \ \ \ \ \ \ \ \ \ \ \ \ \ \ \ \ \ \left(
\begin{array}
[c]{c}%
\text{here, we have renamed the summation index }j\\
\text{as }i\text{ in the second sum}%
\end{array}
\right) \\
&  =\underbrace{\sum_{\left(  i,j\right)  \in\mathbb{N}\times\mathbb{N}}%
}_{=\sum_{n\in\mathbb{N}}\ \ \sum_{\substack{\left(  i,j\right)  \in
\mathbb{N}\times\mathbb{N};\\pj+i=n}}}\left(  -1\right)  ^{j}\dbinom{k}%
{j}\dbinom{i+k-1}{i}x^{pj+i}\\
&  =\sum_{n\in\mathbb{N}}\ \ \sum_{\substack{\left(  i,j\right)  \in
\mathbb{N}\times\mathbb{N};\\pj+i=n}}\left(  -1\right)  ^{j}\dbinom{k}%
{j}\dbinom{i+k-1}{i}\underbrace{x^{pj+i}}_{\substack{=x^{n}\\\text{(since
}pj+i=n\text{)}}}\\
&  =\sum_{n\in\mathbb{N}}\ \ \sum_{\substack{\left(  i,j\right)  \in
\mathbb{N}\times\mathbb{N};\\pj+i=n}}\left(  -1\right)  ^{j}\dbinom{k}%
{j}\dbinom{i+k-1}{i}x^{n}\\
&  =\sum_{n\in\mathbb{N}}\left(  \sum_{\substack{\left(  i,j\right)
\in\mathbb{N}\times\mathbb{N};\\pj+i=n}}\left(  -1\right)  ^{j}\dbinom{k}%
{j}\dbinom{i+k-1}{i}\right)  x^{n}.
\end{align*}
Thus, (\ref{pf.thm.fps.comps.num-wpcomps-n-k.3}) becomes
\[
W_{k,p}=\left(  1-x^{p}\right)  ^{k}\left(  1-x\right)  ^{-k}=\sum
_{n\in\mathbb{N}}\left(  \sum_{\substack{\left(  i,j\right)  \in
\mathbb{N}\times\mathbb{N};\\pj+i=n}}\left(  -1\right)  ^{j}\dbinom{k}%
{j}\dbinom{i+k-1}{i}\right)  x^{n}.
\]
Comparing coefficients in this equality, we find that each $n\in\mathbb{N}$
satisfies%
\begin{align*}
w_{n,k,p}  &  =\sum_{\substack{\left(  i,j\right)  \in\mathbb{N}%
\times\mathbb{N};\\pj+i=n}}\left(  -1\right)  ^{j}\dbinom{k}{j}\dbinom
{i+k-1}{i}\\
&  =\sum_{\substack{j\in\mathbb{N};\\pj\leq n}}\left(  -1\right)  ^{j}%
\dbinom{k}{j}\dbinom{n-pj+k-1}{n-pj}\\
&  \ \ \ \ \ \ \ \ \ \ \ \ \ \ \ \ \ \ \ \ \left(
\begin{array}
[c]{c}%
\text{here, we have substituted }\left(  n-pj,j\right)  \text{ for }\left(
i,j\right)  \text{ in the sum,}\\
\text{since any pair }\left(  i,j\right)  \in\mathbb{N}\times\mathbb{N}\text{
satisfying }pj+i=n\\
\text{is uniquely determined by its second entry }j
\end{array}
\right) \\
&  =\sum_{j\in\mathbb{N}}\left(  -1\right)  ^{j}\dbinom{k}{j}\dbinom
{n-pj+k-1}{n-pj}\\
&  \ \ \ \ \ \ \ \ \ \ \ \ \ \ \ \ \ \ \ \ \left(
\begin{array}
[c]{c}%
\text{here, we have extended the range of summation by}\\
\text{dropping the \textquotedblleft}pj\leq n\text{\textquotedblright%
\ requirement; this does not change}\\
\text{the sum, since all newly introduced addends are }0
\end{array}
\right) \\
&  =\sum_{j=0}^{k}\left(  -1\right)  ^{j}\dbinom{k}{j}\dbinom{n-pj+k-1}{n-pj}%
\end{align*}
(here, we have removed all addends with $j>k$ from the sum; this does not
change the sum, since all these addends are $0$).
\end{proof}

Thus, we have proved the following fact:

\begin{theorem}
\label{thm.fps.comps.num-wpcomps-n-k}Let $n,k,p\in\mathbb{N}$. Then, the \# of
$k$-tuples $\left(  \alpha_{1},\alpha_{2},\ldots,\alpha_{k}\right)
\in\left\{  0,1,\ldots,p-1\right\}  ^{k}$ satisfying $\alpha_{1}+\alpha
_{2}+\cdots+\alpha_{k}=n$ is%
\[
\sum_{j=0}^{k}\left(  -1\right)  ^{j}\dbinom{k}{j}\dbinom{n-pj+k-1}{n-pj}.
\]

\end{theorem}

In general, this expression is the simplest we can get. A combinatorial proof
of Theorem \ref{thm.fps.comps.num-wpcomps-n-k} can be found in \cite[Exercise
2.10.6]{19fco}.

However, the particular case when $p=2$ is worth exploring, as it allows for a
much simpler expression. Indeed, the $k$-tuples $\left(  \alpha_{1},\alpha
_{2},\ldots,\alpha_{k}\right)  \in\left\{  0,1\right\}  ^{k}$ are just the
\textquotedblleft\emph{binary }$k$\emph{-strings}\textquotedblright, i.e., the
$k$-tuples formed of $0$s and $1$s. Imposing the condition $\alpha_{1}%
+\alpha_{2}+\cdots+\alpha_{k}=n$ on such a $k$-tuple is tantamount to
requiring that it contain exactly $n$ many $1$s. Therefore, the \# of all
$k$-tuples $\left(  \alpha_{1},\alpha_{2},\ldots,\alpha_{k}\right)
\in\left\{  0,1\right\}  ^{k}$ satisfying $\alpha_{1}+\alpha_{2}+\cdots
+\alpha_{k}=n$ is $\dbinom{k}{n}$, since we have to choose which $n$ of its
$k$ positions will be occupied by $1$s (and then all remaining $k-n$ positions
will be occupied by $0$s). However, Theorem
\ref{thm.fps.comps.num-wpcomps-n-k} (applied to $p=2$) yields that this \#
equals $\sum_{j=0}^{k}\left(  -1\right)  ^{j}\dbinom{k}{j}\dbinom
{n-2j+k-1}{n-2j}$. Comparing these two results, we obtain the following identity:

\begin{proposition}
\label{prop.fps.comps.num-w2comps-n-k-id}Let $n,k\in\mathbb{N}$. Then,%
\[
\dbinom{k}{n}=\sum_{j=0}^{k}\left(  -1\right)  ^{j}\dbinom{k}{j}%
\dbinom{n-2j+k-1}{n-2j}.
\]

\end{proposition}
